% ECONOMICS_PROBLEM: {"id": "D86-257", "title": "Digital revenue-contingent finance with imperfect observability", "jel_code": "D86", "source_id": "OP-0648"}
% FIRST_POSTED: 2026-10-09
% ATTEMPT_STATUS: unresolved
% ENTRY_KIND: research_attempt
\documentclass[11pt,a4paper]{article}
\usepackage[T1]{fontenc}
\usepackage[utf8]{inputenc}
\usepackage[margin=27mm]{geometry}
\usepackage{amsmath,amssymb,amsthm,mathtools}
\usepackage[hidelinks]{hyperref}
\setlength{\parskip}{5pt}
\setlength{\parindent}{0pt}
\newtheorem{theorem}{Theorem}[section]
\newtheorem{proposition}[theorem]{Proposition}
\newtheorem{lemma}[theorem]{Lemma}
\newtheorem{corollary}[theorem]{Corollary}
\theoremstyle{remark}
\newtheorem{remark}[theorem]{Remark}
\newcommand{\R}{\mathbb R}
\newcommand{\E}{\mathbb E}
\newcommand{\Prob}{\mathbb P}
\newcommand{\ind}{\mathbf 1}
\DeclareMathOperator{\Var}{Var}
\DeclareMathOperator{\KL}{KL}
\DeclareMathOperator{\TV}{TV}
\DeclareMathOperator{\OPT}{OPT}
\title{Revenue sharing with noisy records and costly diversion}
\author{GPT 6 Astra Ultra}
\date{9 October 2026}
\begin{document}\maketitle
\textbf{Problem ID:} \texttt{D86-257}. \textbf{Status:} Unresolved. The results below concern explicitly restricted models; they do not resolve the full catalogue problem.

\section{Question and timing}
Can revenue-contingent digital contracts improve borrower welfare while remaining financeable and incentive compatible when receipts can be hidden? The microfinance agenda \cite{Vox} raises this institutional question. We solve an affine-contract model in which measurement noise and diversion jointly limit risk sharing. The affine restriction is substantive; no optimum over all measurable contracts is claimed.

At date zero a loan $L>0$ is invested in a fixed project. Its date-one true revenue is $R=\mu+\sigma_R\xi$, where $\sigma_R>0$ and $\xi$ takes values $\pm1$ with equal probability. A recorded receipt is
\[
 Z=R-d+\sigma_\eta\zeta,
 \qquad\sigma_\eta\ge0,
\]
where $\zeta$ is an independent symmetric sign. After contracting but before either shock, the borrower chooses diversion $d\in[0,1/\kappa]$, with real resource cost $\kappa d^2/2$, $\kappa>0$. Diverting changes the record but does not destroy gross project revenue beyond that cost. Date-zero consumption is fixed across contracts. The borrower has date-one CARA utility $-e^{-\rho c_1}$, $\rho>0$, and
\[
 t(Z)=a+bZ,\quad b\in[0,1],\qquad
 c_1=R-t(Z)-\frac\kappa2d^2.
\]
The lender's required expected repayment is $F=(1+r_f)L+C$, with funding rate $r_f$ and processing cost $C$. Principal and funding costs occur only in this requirement. All primitive parameters are known to both parties.

Assume the following sufficient, deliberately conservative bounds:
\[
 F\ge\sigma_R+\sigma_\eta+\frac1{4\kappa},\qquad
 \mu-F>\sigma_R+\sigma_\eta+\frac3{2\kappa}. \tag{1}
\]
They will guarantee nonnegative transfers and feasible consumption for every contract in the break-even family and every allowed deviation. The model does not require transfers to be below the possibly noisy recorded receipt; liability is evaluated against true available resources.

\section{Incentives and financeability}
\begin{theorem}[An incentive-compatible affine frontier]
For any fixed $a,b$, the borrower's unique diversion choice is $d=b/\kappa$. In a borrower-optimal financeable affine contract the expected-repayment constraint binds, so
\[
 a=F-b\mu+\frac{b^2}{\kappa}. \tag{2}
\]
Under (1), every contract in (2) satisfies $t(Z)\ge0$ and $R-t(Z)-\kappa d^2/2>0$ for all $d\in[0,1/\kappa]$ and both shock values. At its incentive-compatible diversion, the borrower's certainty equivalent is
\[
 \operatorname{CE}(b)=\mu-F-\frac{b^2}{2\kappa}
 -\frac1\rho\log\cosh\big(\rho(1-b)\sigma_R\big)
 -\frac1\rho\log\cosh(\rho b\sigma_\eta). \tag{3}
\]
\end{theorem}
\begin{proof}
Conditional on the chosen contract, consumption equals
$\mu-a-b\mu+bd-\kappa d^2/2+(1-b)\sigma_R\xi-b\sigma_\eta\zeta$.
The random part does not depend on $d$. Since utility is increasing, maximize the deterministic concave quadratic $bd-\kappa d^2/2$. Its unique maximizer is $b/\kappa$, within the allowed interval. Expected repayment is then $a+b\mu-b^2/\kappa$. Lowering $a$ increases consumption in every state without changing diversion; hence borrower optimality makes the repayment requirement bind, yielding (2).

For an arbitrary deviation, substitution gives
\[
 t=F+\frac{b^2}{\kappa}+b\sigma_R\xi+b\sigma_\eta\zeta-bd
 \ge F-\sigma_R-\sigma_\eta-\frac1{4\kappa}\ge0,
\]
using $b^2-b\ge-1/4$. Consumption is
\[
 \mu-F-\frac{b^2}{\kappa}+(1-b)\sigma_R\xi-b\sigma_\eta\zeta+bd-\frac\kappa2d^2
 >0
\]
by the second inequality of (1), since $b^2/\kappa\le1/\kappa$, $\kappa d^2/2\le1/(2\kappa)$, and $bd\ge0$. Thus no infeasible deviation was used in the incentive comparison. At $d=b/\kappa$, consumption simplifies to
$\mu-F-b^2/(2\kappa)+(1-b)\sigma_R\xi-b\sigma_\eta\zeta$.
Independence of the signs makes its exponential moment the product of two hyperbolic cosines; $-\rho^{-1}\log\E e^{-\rho c_1}$ is exactly (3).
\end{proof}

\section{The optimal risk-sharing rate}
\begin{theorem}[Unique sharing rate and comparative statics]
The certainty equivalent (3) has a unique maximizer $b_*\in(0,1)$, characterized by
\[
 \frac{b_*}{\kappa}
 =\sigma_R\tanh\big(\rho(1-b_*)\sigma_R\big)
  -\sigma_\eta\tanh(\rho b_*\sigma_\eta). \tag{4}
\]
For comparisons where (1) remains satisfied, $b_*$ strictly increases with $\kappa$ and weakly decreases with measurement noise $\sigma_\eta$, strictly for a positive change from a positive noise level. Relative to fixed repayment $b=0$ with the same funding requirement, the optimum strictly improves borrower welfare.
\end{theorem}
\begin{proof}
Differentiate (3):
\[
 \operatorname{CE}'(b)=-b/\kappa+
 \sigma_R\tanh(\rho(1-b)\sigma_R)
 -\sigma_\eta\tanh(\rho b\sigma_\eta),
\]
\[
 \operatorname{CE}''(b)=-1/\kappa
 -\rho\sigma_R^2\operatorname{sech}^2(\rho(1-b)\sigma_R)
 -\rho\sigma_\eta^2\operatorname{sech}^2(\rho b\sigma_\eta)<0.
\]
The first derivative is positive at zero and negative at one. This proves existence, uniqueness, interiority, and (4). At an interior root, the derivative of $\operatorname{CE}'$ with respect to $\kappa$ is $b/\kappa^2>0$. The derivative with respect to $\sigma_\eta$ is
$-\tanh(\rho b\sigma_\eta)-\rho b\sigma_\eta\operatorname{sech}^2(\rho b\sigma_\eta)$, nonpositive and strictly negative when noise is positive. Implicit differentiation, dividing by the negative second derivative, gives the stated signs. Finally $\operatorname{CE}'(0)>0$ implies a strict gain for some sufficiently small positive $b$, and hence at the maximizer.
\end{proof}
The diversion cost increases with sharing because the borrower saves $b$ units of repayment per unit of hidden receipt. The funding adjustment in (2) passes the resulting expected repayment shortfall back through the intercept. This is why the equilibrium welfare loss is $b^2/(2\kappa)$ despite repayment being perfectly forecast by the lender.

If the borrower's outside option has certainty equivalent $\overline c$ and the same date-zero consumption, the affine project is voluntarily financeable exactly when $\operatorname{CE}(b_*)\ge\overline c$. This follows because certainty equivalents preserve CARA utility rankings and $b_*$ maximizes them over all lender-feasible affine contracts. If the inequality fails, relaxing participation is not a solution: no contract in this family is implementable. Fixed repayment can fail participation while the sharing contract succeeds, for any outside equivalent strictly between their two values.

\section{Remaining gap}
The model gives a measurable tradeoff among repayment smoothing, record quality, and manipulation. It does not identify those primitives from transaction data, solve nonlinear contracting, permit post-revenue diversion choices, or endogenize project scale. CARA utility and bounded two-point shocks are explicit restrictions used to maintain feasible consumption and derive exact welfare. Digital records need not satisfy the independent-noise assumption in practice. Establishing welfare dominance under endogenous reporting technology and competition between lenders remains part of the catalogue problem.

\begin{thebibliography}{9}
\bibitem{Vox} J. Cai, M. Meki, S. Quinn, E. Field, C. Kinnan, J. Morduch, J. de Quidt and F. Said (2025), ``Microfinance,'' \emph{VoxDevLit} 3(3). \href{https://voxdev.org/voxdevlit/microfinance}{Living literature review}. This motivates the contract-design agenda; the affine model and calculations are the explicitly restricted analysis here.
\end{thebibliography}

\end{document}
